\chapter{Outdoor Exposure and Oscillating Temperatures}

A time honored procedure in horticulture and gardening has been to sow seeds
in pots and place them outdoors in the fall. Germination often takes place in the
spring. In light of the discussions up to here it is easy to see why this procedure is so
often successful. It provides the low temperatures at which some species germinate
and the low temperatures where some inhibitor systems are destroyed. It is also
moderately successful with seeds with impervious seed coats as described in the
previous chapter. However, there is much more to it.

A number of species were found which gave good germination when placed
outdoors whereas \textit{germination was low or none in all other treatments}. The following
list gives some idea of the variety of species that showed this effect: Acer
pennsylvanicum, Agalinis setacea, Aquilegia flavescens, Aquilegia jucunda, Aquilegia
olympica, Arisaema dracontium, Arisaema quinquifolium, Arisaema triphyllum, Asarum
canadensis, Collinsoniacariadensis, Cornus alternifolia, Cornus amomum, Cornus
sibirica, Echinocystis lobata, Fritillaria pudica, Juglans cinerea, Juglans nigra,
Impatiens biflora, Iris forrestii, Iris lactea, Iris setosa, Linum capitatum Myrrhis odorata,
Phellodendron amurense, Phlox paniculata, Sanguinaria canadensis, Trollium laxus,
and Xerophyllum asphodeloides.

Now there is nothing magical about being outdoors, and there must be some
explanation in terms of chemical rate theory. One explanation is that the oscillating
temperatures themselves are required. This explanation has been proposed before
for the germination of Agrostis alba and Lycopus europaeus (ref. 16) based on limited
data. The other explanation is that germination occurs at some temperature
intermediate between 40 and 70. The oscillating temperatures pass through this
intermediate temperature and provide enough residence time at the intermediate
temperature to achieve germination. The results described below for Cleome
serrulata favor the first explanation, namely that the oscillating temperatures
themselves are required. The theory of this is discussed in Chapter 19.

It had already been found with Cleome serrulata that seeds collected in
September and placed outdoors germinated in October whereas controls held at
either 40 or 70 failed to germinate a single seed. Specifically, seeds were collected
on September 27, 1990. Portions were placed in moist towels at 70, 40, and in
outdoor conditions. On October 12, fifteen days later 125/125 (100\%) of the seeds
placed outdoor had germinated vigorously. In this period outdoor temperatures had
ranged from 30 to 70 with large daily oscillations. There was absolutely no
germination in either 70 or 40 in the dark. The experiment was repeated on a sample
of seed collected October 23, 1990. The outdoor treatment germinated 36\% on
November 11, 7\% on November 15, and 11\% on December 3. At that time
temperatures dropped below freezing each day and germination ceased. Again there
was not a single germination in the samples in the dark at 70 and 40.

The experiment was now repeated in a more controlled manner using exactly
controlled oscillating temperatures. The Cleome seeds were subjected to alternating
twelve hours at 70 and twelve hours at 40 starting one sample at 40 and the other at
70. Both samples of seeds germinated following a zero order rate law with a rate of
6\% per day and an induction time of four days.

The experiment was now made more elegant by subjecting the seeds to six
hour, twelve hour, and twenty four hour periods alternating between temperatures of
70 and 40 and starting at either 40 or 70. Germination started after four days in all six
experiments. Although the rate was somewhat faster at the start for the six hour
alternations, at the end of a week the rate curves were virtually coincident. The
experiment in which the temperature shifts were every six hours spent more time at
intermediate temperatures. Since its rate was about the same as the others, it
indicates that the oscillating temperatures themselves are responsible for the
germination. Obviously experiments should be conducted at constant temperatures of
40, 50 ,60, and 70. These are planned, and these will decisively settle the question.

Germination of the .Cleome seeds is also promoted by light as described in Chapter
20. Why light and oscillating temperatures should have comparable efficacy in
initiating germination is a chemical mystery that will someday be solved.
With most examples of germination in outdoor treatment, the seeds were placed
outdoors in the fall with germination occurring in spring. However, with a few such as
Cleome serrulata and Trillium flexipes germination occurred in September or October
before anything like winter conditions had begun. It was the results with these two
species that indicated that the temperature oscillations themselveswere the key factor.
Where seeds were placed outdoors in early fall but did not germinate until spring, the
interpretation is that a chilling period around 40 was needed first in order to destroy
germination inhibitors followed by the period of oscillating temperatures.
Corydalis nobile has an interesting pattern. The seeds do not tolerate dry
storage so they must be placed in the moist towels when they ripen in June. In the
oscillating temperatures of fall about half of them expand enought to split the seed
coats. Only these split seeds germinate in the oscillating temperatures of the following
spring. The oscillating temperatures of both fall and spring are required.
With the discovery of the effect of oscillating temperatures, outdoor treatment
became one of the standard procedures. It was found that many of the species that
had previously given low and much extended germinations gave high germination in
the outdoor treatment.

